% Generated draft: requires mathematical and source review.
\documentclass[11pt]{article}
\usepackage[margin=25mm]{geometry}
\usepackage{fontspec}
\setmainfont{Latin Modern Roman}
\usepackage{amsmath,amssymb,mathtools,unicode-math}
\setmathfont{Latin Modern Math}
\usepackage{enumitem,needspace,xurl,hyperref}
\setlength{\parindent}{0pt}
\setlength{\parskip}{5pt}
\setlength{\emergencystretch}{3em}
\newcommand{\N}{\mathbb{N}}
\newcommand{\Z}{\mathbb{Z}}
\newcommand{\Q}{\mathbb{Q}}
\newcommand{\R}{\mathbb{R}}
\newcommand{\C}{\mathbb{C}}
\newcommand{\F}{\mathbb{F}}
\newcommand{\A}{\mathbb{A}}
\newcommand{\PP}{\mathbb{P}}
\newcommand{\E}{\mathbb{E}}
\newcommand{\eps}{\varepsilon}
\newcommand{\dd}{\,\mathrm{d}}
\begin{document}
\section*{UnsolvedMath ID: 87}
\textbf{Source problem code:} catalogue entry
\subsection{Definitions and mathematical statement}
A number field is a finite field extension $K/\Q$. For an integer $n\geq1$, a quadratic form in $n$ variables over $K$ is a homogeneous polynomial
\[
 q(x_1,\ldots,x_n)=\sum_{i=1}^{n}a_{ii}x_i^2+
 \sum_{1\leq i<j\leq n}a_{ij}x_ix_j,\qquad a_{ij}\in K.
\]
Since $K$ has characteristic zero, one may equivalently write $q(x)=x^{\mathsf T}Ax$ with $A$ a symmetric matrix over $K$. A finite set of algebraic numerical coefficients lies in some number field, so this framework expresses the coefficient condition in the supplied statement without restricting to rational coefficients.

For clarification of the arithmetic terminology, two forms $q,q'$ in $n$ variables are equivalent over $K$ if there is an invertible matrix $T\in\operatorname{GL}_n(K)$ such that $q'(x)=q(Tx)$ for all $x\in K^n$. A form $q$ represents $b\in K$ over $K$ if there exists $x\in K^n$ with $q(x)=b$. These definitions illustrate classification and representation questions; the source does not select a particular equivalence problem or represented element.

\textbf{Historical programme (Hilbert, 1900).} Extend the arithmetic theory of quadratic forms to forms whose coefficients are algebraic numbers, equivalently to forms over arbitrary number fields as above. The supplied record identifies class field theory and the Hasse--Minkowski theorem as relevant historical progress, but gives no particular residual representation assertion to prove. It supplies no nondegeneracy, integrality, signature, dimension, or prescribed-representation restrictions. Such restrictions must not be inferred as part of one unstated conjecture; integral and field-valued representation problems in particular have different hypotheses and conclusions.

The imported background says that general representation questions remain open, whereas the literature triage dated 2026-08-17 classifies the classical programme as \emph{solved in the literature}, conventionally resolved through the mature arithmetic theory and local--global classification tools. The triage explicitly says that later representation questions do not reopen the original programme and that this import specifies no single residual claim. The registry still labels the item \emph{open}. These conflicting source-level status descriptions are recorded without choosing a new current status or turning the programme into an unsupported open conjecture.
\subsection{Short English statement}
Hilbert's eleventh problem asks for an arithmetic theory of quadratic forms with algebraic-number coefficients. The record is a broad historical programme rather than a particular representation conjecture. Its dated review treats the classical programme as resolved and distinguishes later questions, despite the imported open label.
\subsection{Sources}
\begin{itemize}
\item UnsolvedMath local registry, data/unsolvedmath/problems.json, permanent ID 87, source code catalogue entry, ``Hilbert's 11th Problem: Quadratic Forms over Algebraic Number Fields,'' attributed to David Hilbert (1900).
\item Encyclopedia of Mathematics, ``Hilbert problems,'' \url{https://encyclopediaofmath.org/wiki/Hilbert_problems}. The supplied review records access on 2026-08-17 and says its status overview lists Problem 11 among consensus resolutions. This draft has not independently consulted that reference.
\item Ulam AI and the UnsolvedMath Contributors, \emph{UnsolvedMath: A Curated Collection of Open Mathematics Problems}; this local record includes a literature triage dated 2026-08-17. Dataset curation and original metadata are supplied under CC BY 4.0; referenced source material retains its own terms.
\end{itemize}
\end{document}
